\mainlanguage[de]

\setuplayout
  [backspace=2.0cm,
   height=middle,
   width=middle,
   leftmargin=1.8cm,
   footerdistance=0.0cm,
   rightmargindistance=0.2cm,
   rightmargin=1.8cm,
   bottom=0.5cm,
   leftmargindistance=0.2cm,
   header=2.1cm,
   headerdistance=0.5cm,
   footer=0.5cm,
   top=0cm,
   topspace=0.5cm,
   bottomspace=0.0cm]

\setupformulas[align=right]

\definehead[exercise][subsection]

\setuphead
  [exercise]
  [number=yes,
   style=,
   numberstyle=\ss\bf\tfb,
   textstyle=\ss,
   alternative=margin,
   numbercommand={\determineheadnumber[subsection]\currentheadnumber\gobbleoneargument}]

\setuphead
  [chapter]
  [number=no,
   style=,
   textstyle=\tfb\ss]

\setuphead
  [section]
  [number=no,
   style=,
   textstyle=\tf\ss]

\definemeasure[headerimageheight][\dimexpr(\headerheight-2ex)]

\definelayer[sabel]

\startsetups[sabel]

\setlayerframed[sabel][x=0mm,y=0mm][frame=off,offset=none]
{%
   \setupTABLE[start][style=\tfx]
   \setupTABLE[r][each][height=0.85\baselineskip]
   \setupTABLE[column][1][align=right, width=.4\textwidth, loffset=-2pt, toffset=-5pt]%
   \setupTABLE[column][2][align={center,lohi}, width=.2\textwidth]%
   \setupTABLE[column][3][align=left, width=.4\textwidth, roffset=-2pt, toffset=-5pt]%
   \bTABLE[frame=off]
   \bTR \bTD Sabel-GBS       \eTD \bTD [nr=5] \dontleavehmode\externalfigure[assets/sabel.png][maxheight=\measure{headerimageheight}, frame=off]\dontleavehmode  \eTD \bTD Semester WS 2025 \eTD \eTR
   \bTR \bTD Messtechnik     \eTD                                                                                                                                      \bTD Klasse E1A       \eTD \eTR
   \bTR \bTD B.~Donner       \eTD                                                                                                                                      \bTD                  \eTD \eTR
   \bTR \bTD \strut          \eTD                                                                                                                                      \bTD                  \eTD \eTR
   \bTR \bTD 1. Schulaufgabe \eTD                                                                                                                                      \bTD 20.11.2025       \eTD \eTR
   \eTABLE}

\stopsetups

\setupbackgrounds[header][text][bottomframe=on,setups=sabel,background=sabel]
\setuppagenumbering[location=]

\setupfootertexts
  [BE: 31 Arbeitszeit: 45min][Viel Erfolg!]
  [Viel Erfolg!][BE: 31 Arbeitszeit: 45min]


\definedescription[example][
  headstyle=\it, style=, align=left, location=top,
  width=broad, margin=1cm, indenting=never, text=\hskip-1cm]

% \definemathmatrix[pmatrix][
%   left={\left(\,}, right={\,\right)}]

%\usemodule[units]

\definemathcommand [kgV] [nolop] {\mfunction{kg\kern-.15em\hbox{V}}}

\setdigitmode 3

\setupprefixtext[en][micro=μ]

\starttext

\exercise{Analoge und digitale Messgeräte}

Ein analoges Voltmeter soll hinsichtlich der Genauigkeit mit einem digitalen
Multimeter verglichen werden.
Der Vergleich soll für beide Geräte in {200}\unit{micro Volt} Messbereich
durchgeführt werden.
Das analoge Gerät besitzt die Genauigkeitsklasse {\digits 0.5 }. Das
digitale Gerät hat ein \math{3\frac{1}{2}}~stelliges Display und einem Fehler
von \math{\pm(2 \%\text{v. MW} + 2 \text{Digits})}:

\startitemize[a]
\item Bei welchem Messwert haben beide Instrumente den gleichen absoluten
Fehler (mit Herleitung!)? (BE~4)
\item Begründen Sie, welche relativen Fehler sich bei diesem Messwert für das
analoge und das digitale Gerät ergeben. (BE~2)
\item Welches von beiden Geräten hat einen kleineren absoluten Fehler, wenn
der unter a) gefundene Werte überschritten wird? Begründen Sie Ihre
Antwort. (BE~2)
\item Welches von beiden Geräten hat einen kleineren relativen Fehler, wenn der
unter a) gefundene Werte überschritten wird? Begründen Sie Ihre Antwort. (BE~2)
\stopitemize

\exercise{Widerstandsmessung}

Ein unbekannter Widerstand \math{R} soll durch eine Strom- und Spannungsmessung
ermittelt werden. Verwendet werden ein Spannungsmesser mit dem Innenwiderstand
\math{60 \unit{kilo ohm}} und ein Strommesser mit Innenwiderstand
{\digits 300 }\unit{ohm}. Es wird eine Spannung von
{\digits 15.3 }\unit{Volt} und ein Strom {\digits 1.65 }\unit{mA}
gemessen. 

\startitemize[a]
\item Zeichnen Sie die spannungsrichtige {\bf und} stromrichtige Schaltung mit
den jeweiligen Ersatzschaltbildern der verwendeten Messgeräte und beschrifte
die Zeichung. (BE~4)
\item Wie groß wäre der naive Widerstandswert von \math{R}? (BE~1)
\item Wie groß wäre der wahre Widerstandswert von \math{R}, wenn diese Werte
in der stromrichtigen Schaltung gemessen werden würden? (BE~1)
\item Wie groß wäre der wahre Widerstandswert von \math{R}, wenn diese Werte
in der spannungsrichtigen Schaltung gemessen werden würden? (BE~1)
\stopitemize

% \exercise{Widerstandsbestimmung}

% In der folgenden Aufgabe soll ein Widerstand mit ungefähr
% {\digits 100 }\unit{ohm} möglichst genau vermessen werden. Ihnen stehen zwei
% Messgeräte zur Verfügung, wobei das Spannungsmessgerät einen Innenwiederstand
% von {\digits 10 }\unit{Mega ohm}, das Strommessgerät hat in dem fraglichen
% Messbereich einen Innenwiderstand von {\digits 1 }\unit{ohm}.\\
% Können Sie den Widerstand genauer mit einer spannungsrichtigen oder einer
% stromrichtigen Schaltung vermessen. Schätzen Sie die zu erwartenden Fehler in
% beiden Schaltungen ab und begründen Sie damit Ihre Entscheidung. (BE~3)

\page[yes]

\exercise{Reihenschaltung von Widerständen}

Ein Widerstand \math{R_1 = \text{{\digits 900 }\unit{ohm}}} mit einer relativen
Fehlergrenze von 3\% wird mit einem weiteren Widerstand
\math{R_2 = \text{{\digits 1.1 }\unit{kilo ohm}}} mit einer relativen
Fehlergrenze von
6\% in Reihe geschaltet.

\startitemize[a]
\item Wie groß ist der Gesamtwiderstand und dessen relative
Fehlergrenze? (BE~4)
\item Welche relative Fehlergrenze bräuchte der Widerstand \math{R_2}, wenn
\math{R_1}, \math{R_2} und der relative Fehler von \math{R_1} unverändert
bleiben, damit der relative Fehler des Gesamtwiderstands 2\% beträgt? (BE~3)
\stopitemize

\exercise{Leistungsmessung}

Der Spannungsmesser hat einen Innenwiderstand \math{50 \unit{kilo ohm}} mit
relativer Fehlergrenze von {\digits 0.01 }\% und der Strommesser einen
Innenwiderstand von {\digits 15 }\unit{ohm} und relativer Fehlergrenze von
{\digits 0.5 }\%. Es wird eine Spannung von {\digits 3.5 }\unit{Volt} mit
relativer Fehlergrenze von 5\% und ein Strom {\digits 235 }\unit{mA} mit
relativer Fehlergrenze von 9\% gemessen. 

\startitemize[a]
\item Zeigen Sie durch eine sorgfältige Herleitung, dass für die korrekte
Leistung in der spannungsrichtigen Schaltung gilt:\\
\math{P = U\cdot I - \frac{U^2}{R_{IV}}}. (BE~3)
\item Geben Sie den absoluten Fehler auf 6 Nachkommastellen an, der sich für
die korrekte Leistung nach der oben genannten Beziehung ergeben würde. (BE~4)
%\item Welcher absoluter Fehler würde sich für die Leistung ergeben wenn man die Leistung nach der Formel \math{P = U^2\left(\frac{I}{U} - \frac{1}{R_{IV}}\right)} bestimmt? Ist die Beziehung falsch? Woher rührt der Unterschied? (BE~4)
\stopitemize

\stoptext
